\documentstyle[%


righttag%


,11pt%


,amssymb%


,russian%               remove in Eng.


,izvru%                 Set izven% in Eng.


]{amsart}





% 





\newcommand{\gea}{\geqslant}


\newcommand{\lea}{\leqslant}


\newcommand{\vp}{\varphi}


\newcommand{\vk}{\varkappa}


\newcommand{\ve}{\varepsilon}


\newcommand{\si}{\sigma}


\newcommand{\N}{\Bbb{N}}


\newcommand{\R}{\Bbb{R}}


\newcommand{\bC}{\Bbb{C}}


\newcommand{\bH}{\Bbb{H}}


\newcommand{\rL}{\cal{L}}


\newcommand{\rB}{\cal{B}}


\newcommand{\rP}{\cal{P}}





\newcommand{\ctg}{\operatorname{ctg}}


\newcommand{\tts}[1]{}


\renewcommand{\baselinestretch}{1.1}





%\hoffset=-15mm%                  For Eng.


\setcounter{page}{39}


\god{2004}


\nomer{8~(507)} %                        Remove in Eng.


%\nomer{Vol.\,48, No.\,8}%               Set in Eng.


%\str{\pageref{begin}--\pageref{end}}%  Set in Eng.


\udk{517.968:519.6}





\author{\it Л.Б.\,Ермолаева}


% NB:   ^^ remove in Eng.


\title{Решение периодическиx краевых задач методом подобластей}





\date{}


%\pagestyle{myheadings}%         Set in Eng.


\pagestyle{plain} %              Remove in Eng.


\begin{document}


%\thispagestyle{plain}%          Set in Eng.


\maketitle


\label{begin}








Данная заметка, являющаяся продолжением [1], [2], посвящена теоретическому


обоснованию метода подобластей решения периодических краевых задач для


операторно-диф\-фе\-рен\-ци\-аль\-ных уравнений с 


периодическими коэффициентами. При этом


основное внимание уделено доказательству сходимости метода в пространствах


С.Л.\,Соболева и установлению эффективных (в том числе неулучшаемых по порядку)


оценок погрешности в ряде функциональных пространств, учитывающих структурные


свойства коэффициентов исследуемых уравнений.





Рассмотрим периодическую краевую задачу для уравнения


\begin{gather}


Ax \equiv x^{(m)}(s) + Bx(s)=y(s), \quad - \infty < s < \infty,\\


%


x^{(k)}(0)=x^{(k)} (2 \pi), \quad k=\overline{0,m-1},


\end{gather}


где $y \in L_{1}=L_{1}(0,2 \pi)$, $B$ --- некоторый линейный (напр.,


интегродифференциальный) оператор с областью значений в пространстве


$L_{1}$, а $m$ --- целое неотрицательное число, причем при $m=0$ условия


(2) отсутствуют.





Приближенное решение задачи (1)--(2) ищем в виде


полинома


\begin{equation}


x_n(s)=\frac{a_{0}}{2} + \sum_{k=1}^n a_k \cos ks + b_k \sin ks


\equiv \sum_{k=-n}^n c_k e^{iks},


\end{equation}


коэффициенты которого определим из условий


\begin{equation}


\int_{s_j}^{s_{j+1}} (Ax_n - y)(s) ds=0, \quad j=\overline{0, 2n},


\end{equation}


где


\begin{equation}


s_k=\frac{2k \pi}{2n+1}, \quad k=\overline{0, 2n}, \ \ n \in N.


\end{equation}





Ясно, что условия (4), (5) эквивалентны системе линейных алгебраических


уравнений (СЛАУ) порядка $ 2n+1 $ относительно коэффициентов полинома (3):


\begin{equation}


\sum_{k=-n}^n \alpha_{jk} c_k=y_j, \quad j=\overline{0, 2n},


\end{equation}


где


$$


\alpha_{jk}=\int_{s_j}^{s_{j+1}} (Ae^{iks})(s) ds, \quad y_j=\int_{s_j}^


{s_{j+1}} y(s) ds.


$$








\begin{thm}


Пусть выполнены условия


\begin{itemize}


\item[а)] $y(s) \in L_2$, а оператор $B:W_2^m \to L_2$ вполне непрерывен$;$


\item[б)] задача $(1)$--$(2)$ при $y(s) \equiv 0$ имеет лишь тривиальное решение.


\end{itemize}





Тогда при всех $n$, начиная с некоторого, СЛАУ $(6)$ имеет единственное


решение $\lbrace c_k^* \rbrace_{-n}^n$. Приближенные решения $x_n^*


(s)$ $($т.\,е. $(3)$ при $c_k=c_k^*$, $k=\overline{-n,n})$ сходятся в


пространстве $W_2^m$ к точному решению $x^*(s)$ задачи $(1)$--$(2)$ со


скоростью


\begin{equation}


\| x^* - x_n^* \|_{W_2^m}=O \lbrace E_n^T(x^{*(m)})_2 \rbrace,


\end{equation}


где $E_n^T( \phi)_2$ --- наилучшее среднеквадратическое приближение функции


$\phi(s) \in L_2$ тригонометрическими полиномами порядка не выше $n$.


\end{thm}








\begin{pf}


Введем функциональные пространства Лебега $Y=L_2(0,2 \pi)\equiv


L_2$ и Соболева $X=W_2^m(0,2 \pi)\equiv W_2^m$ с нормами соответственно


\begin{gather*}


\| y \|_Y=\left\{\frac{1}{2 \pi} \int_0^{2 \pi} |y(s)|


^2 ds\right\}^{1/2} \equiv \| y \|_2, \quad y \in L_2;\\


%


\| x \|_X=\sum_{i=0}^m \| x^{(i)}(s) \|_2,


\quad x \in W_2^m \ \ (W_2^0 \equiv L_2),


\end{gather*}


где элементы $x \in X$ удовлетворяют условиям (2). Тогда задачу (1)--(2) можно


записать в виде эквивалентного ей линейного операторного уравнения


\begin{equation}


Ax \equiv Ux + Bx=y \quad (x \in X,\ \ y \in Y),


\end{equation}


где $Ux=x^{(m)}$ при условиях (2). Обозначим через $X_n$ и $Y_n$


подпространства тригонометрических полиномов вида (3) с нормами пространств


$X$ и $Y$ соответственно. Пусть $P_n:L_2 \to Y_{n} \subset L_2$ ---


оператор метода подобластей по тригонометрической системе функций с узлами (5).


Тогда $P_n Ux_n=Ux_n$ для любых $x_n \in X_n$, и поэтому условия (4),


а потому и СЛАУ (6), эквивалентны линейному операторному уравнению


\begin{equation}


A_nx_n \equiv Ux_n + P_nBx_n=P_ny \quad (x_n \in X_n,\ \ P_ny \in Y_n).


\end{equation}





В силу условий теоремы оператор $A:W_2^m \to L_2$ линейно обратим и


обратный оператор $A^{-1}:L_2 \to W_2^m$ ограничен. Покажем, что


операторы $A_n:X_n \to Y_n$ также линейно обратимы (хотя бы при всех


достаточно больших $n$), а обратные операторы $A_n^{-1}:Y_n \to X_n$


ограничены по норме в совокупности:


\begin{equation}


\| A_n^{-1} \|=O(1).


\end{equation}





С этой целью из (8) и (9) для $\forall x_n \in X_n$, $x_n \neq 0$, находим


\begin{multline}


\| Ax_n - A_nx_n \|_Y=\| Bx_n - P_nBx_n


\|_Y \leq \| x_n \|_X \sup_{x \in X,\; \| x \|_X=1}


\| Bx-P_nBx \|_Y=\\


%


=\alpha_n \| x_n \|_X, \quad \alpha_n \equiv \sup


\lbrace \| z-P_nz \|_Y:z \in BS(0,1) \rbrace,


\end{multline}


где $S(0,1)$ -- единичный шар пространства $X$. Будем существенно


опираться на следующие неравенства ([1]; [2], гл.\,1, \S\,1):


\begin{gather}


E_n^T( \varphi)_2 \leq \| \varphi - P_n \varphi \|_2 \leq \frac{\pi}{2} E_n^T


( \varphi)_2, \quad \varphi \in L_2,\quad n\in\N; \\


%


1 \leq {\| P_n \|}_{L_2 \to L_2} \equiv {\| P_n \|}_2 < \frac


{\pi}{2}, \quad n\in\N.


\end{gather}





В силу (12), (13), условия а) теоремы 1 и известной теоремы И.М.\,Гельфанда


(напр., [3], с.\,274--275) в (11) имеем


$\alpha_n \to 0$ при $n \to \infty$.


Поэтому для всех $n \geq n_0$ выполняется неравенство


$$


q_n=\| A^{-1} \| \alpha_n < 1/2 ,


$$


а тогда для таких $n$ в силу теоремы 7 ([4], гл.\,1) операторы $A_n$ линейно


обратимы и


$$


\| A_n^{-1} \| \leq \| A^{-1} \| (1-q_n)^{-1}


\leq 2 \| A^{-1} \|, \quad n \geq n_{0},


$$


т.\,е. соотношение (10) доказано. Поэтому в силу теоремы 6 ([4], гл.\,1) имеем


$$


\| x^* - x_n^* \|_X=\| A^{-1}y - A_n^{-1}P_ny


\|_X \leq \| E-A_n^{-1}P_nA \|\, \| x^* - \ov x_n \|_X,


$$


где $\ov x_n \in X_n$ --- произвольный элемент. Отсюда с учетом


(10)--(13) находим


\begin{equation}


\| x^* - x_n^* \|_X=O \lbrace \| x^* - \ov x_n\|_X \rbrace.


\end{equation}





Для любой функции $x(s) \in W_2^m$ легко доказываются следующие неравенства:


\begin{equation}


\| x^{(k)}(s) \|_{2} \leq \| x^{(k)}(s) \|_{C} \leq


d_0 \| x^{(m)}(s) \|_{2}, \quad k=\overline{0,m-1}, \ \ m \geq 1,


\end{equation}


где $d_0$ --- положительная постоянная, не зависящая от $x(s)$.


Из (14) и (15) получаем оценку


\begin{equation}


\| x^* - x_n^* \|_X=O( \| x^{*(m)}(s) - \ov x_n^{(m)}(s)


\|_{L_2}).


\end{equation}





Выберем $\ov x_n \in X_n$ так, чтобы элемент $\overline x_n^{(m)}


\in Y_n$ был полиномом наилучшего среднеквадратического приближения для


функции $x^{*(m)}\in Y$. Тогда из (16) следует оценка (7).


\end{pf}








\begin{cor}


В условиях теоремы 1 метод (1)--(6) сходится в


пространствах $C_{2 \pi}^k=C^k$ ($C^0=C=C_{2 \pi}$), $k=\overline


{0,m-1}$, со скоростью


\begin{equation}


\| x^* - x_n^* \|_{C^k}=\sum_{i=0}^k \| x^{*(i)} - x_n^{*(i)}


\|_{C}=O \lbrace E_n^T(x^{*(m)})_{2} \rbrace.


\end{equation}


Если же $y \in C_{2 \pi}$, а оператор $B:W_2^m \to C_{2 \pi}$ ограничен,


то справедлива оценка


\begin{equation}


\| x^* - x_n^* \|_{C^m} \equiv \sum_{i=0}^m \| x^{*{(i)}} -


x_n^{*(i)} \|_{C}=O \lbrace E_n^T(x^{*(m)})_{C} \ln n \rbrace,


\end{equation}


где $E_n^T( \varphi)_{C}$ --- наилучшее равномерное приближение функции


$ \varphi(s) \in C $ тригонометрическими полиномами порядка


не выше $n$ ($n+1 \in N$).


\end{cor}








\begin{cor}


Если оператор $B$ и правая часть $y(s)$ уравнения


(1) таковы, что $x^{*(m)}(s) \in H_2^{r+ \alpha }$, то приближенные решения


сходятся в пространствах $W_2^m$, $C^k$, $k=\overline{0,m-1}$,


и $C^m$ со скоростями соответственно


\begin{alignat}{2}


\| x^* - x_n^* \|_{W_2^m} &=O\bigg( \frac{1}{n^{r+ \alpha }}\bigg),


& \quad r+ \alpha &> 0;\\


%


\| x^* - x_n^* \|_{C^k} &=O\bigg( \frac{1}{n^{r+ \alpha}}\bigg),


& \quad k=\overline{0,m-1}, \ \ r + \alpha &> 0;\\


%


\| x^* - x_n^* \|_{C^m} &=O\bigg( \frac{1}{n^{r+ \alpha - 1/2}}\bigg),


& \quad r+\alpha &> \frac{1}{2}.


\end{alignat}


Если же $x^{*(m)}(s) \in H_{\infty}^{r+ \alpha } \equiv W^r H^{\alpha }$,


то метод (1)--(6) сходится в пространстве $C^m$ со скоростью


\begin{equation}


\| x^* - x_n^* \|_{C^m}=O\bigg( \frac{\ln n}{n^{r+ \alpha }}\bigg), \quad


r+ \alpha > 0.


\end{equation}


\end{cor}








{\bf Доказательство следствия 1.}


Из соотношений (7) и (15) получаем оценку (17). Теперь


докажем оценку (18). Из очевидных тождеств


$$


x^{*(m)}(s) \equiv y(s) -(Bx^*)(s),


\quad x_n^{*(m)} \equiv P_ny(s) - (P_nBx_n^*)(s)


$$


и оценки (7) последовательно находим


\begin{multline}


\| x^{*(m)} - x_n^{*(m)} \|_{C} \leq \| x^{*


(m)} - P_nx^{*(m)} \|_{C} + \| P_nB(x^*- x_n^*) \|_{C} \leq \\


%


\leq 2 \| P_n \|_{C}E_n^T (x^{*(m)})_{C} + \| P_n


\|_{C} \| B \|_{W_2^m \to C} \| x^* - x_n^* \|_{W_2^m}=\\


%


=O \lbrace \| P_n \|_{C} ( E_n^T(x^{*(m)})_{


C} + \| x^* - x_n^* \|_{W_2^m}) \rbrace=\\


%


=O \lbrace \| P_n \|_{C}(E_n^T(x^{*(m)})_{C} + E_n^T


(x^{*(m)})_{2}) \rbrace=O \lbrace E_n^T(x^{*(m)})_{C}


\| P_n \|_{C} \rbrace.


\end{multline}





Далее существенно используются неравенства ([2], с.\,41--43)


\begin{gather}


\| S_n \|_{C \to C} \leq \| P_n \|_{C \to C} \leq (1+ \pi)


\| \rL_n \|_{C \to C}, \quad n\in\N; \\


%


\| P_n \|_{C \to C} \equiv \| P_n \|_{C} \asymp \ln n, \quad n \to \infty,


\end{gather}


где для любой $\varphi \in C$


\begin{align*}


S_n(\varphi,s) &=\frac{1}{2 \pi} \int_0^{2 \pi} \frac{\sin \frac{2n+1}{2}


(s- \sigma)}{\sin \frac{s- \sigma}{2}} \varphi( \sigma) d \sigma, \\


%


\rL_n( \varphi,s) &=\frac{1}{2n+1} \sum_{k=0}^{2n} \frac{\sin \frac{2k+1}{2}


(s-s_k)}{\sin \frac{s-s_k}{2}} \varphi(s_k),


\end{align*}


а узлы $s_k$ определены в (5).





Из соотношений (23)--(25) получаем


$$


\| x^{*(m)} - x_n^{*(m)} \|_{C}=O \lbrace E_n^T


(x^{*(m)})_{C} \ln n \rbrace.


$$


Отсюда и из (17) находим оценку (18):


\begin{multline*}


\| x^* - x_n^* \|_{C^m}=\| x^{*(m)} - x_n^{*(m)} \|_{C} +


\| x^* - x_n^* \|_{C^{m-1}}=\\


%


=O \lbrace E_n^{T}(x^{*(m)})_{C} \ln n \rbrace + O \lbrace E_n^T(x^{*(m)})


_{C} \rbrace=O \lbrace E_n^T(x^{*(m)})_{C} \ln n \rbrace.\quad\square


\end{multline*}





{\bf Доказательство следствия 2.} Оценки (19), (20)


следуют из теоремы Джексона в $L_2$ (напр., [5]--[7]) и оценок


соответственно (7), (17). Так как $ H_2^{r + \alpha} \subset H^{r+ \alpha -


1/2}=W^{r_1} H^{\beta}$, где $0 < \beta \leq 1$, $r_{1}=[r + \alpha


- 1/2]$, $r_{1} + \beta=r + \alpha - 1/2$ (напр., [8]), то из теоремы


Джексона в $C_{2 \pi} $ и оценки (18) следует оценка (21). Оценка (22) следует


из (18) и теоремы Джексона в $ C_{2 \pi} $ (напр., [5]--[7]).$\quad\square$








\begin{note}


Доказанная теорема является довольно общей. Здесь в качестве $B$ можно 


взять любой линейный вполне непрерывный оператор из $W_2^m$ в $L_2$, 


например, интегродифференциальный оператор вида


\begin{equation}


Bx=\sum_{k=0}^{m-1} a_k(s) x^{(k)}(s) + \sum_{i=0}^{m} \int_{0}^{2 \pi} h_i


(s, \sigma) x^{(i)}( \sigma) d \sigma,


\tag{$1'$}


\end{equation}


где


$$


h_k(s,\sigma)=h_{0k}(s,\sigma)\ctg\frac{\sigma-s}{2}\ \ (k=\overline


{0,m-1}),\quad h_m(s, \sigma)=h_{0m}(s,\sigma) \bigg| \ctg \frac{\sigma-s}{2}


\bigg|^{\lambda}\ \ (0 \leq \lambda < 1),


$$


причем


$$


a_k(s) \in L_2(0,2 \pi) \ \ (k=\overline{0,m-1}),


\quad h_{0k}(s, \sigma) \in C[0,2\pi]^2\ \ (k=\overline{0,m}).


$$


Здесь возможен также случай $m=0$ (тогда полагаем $\sum\limits_{k=0}^{-1} \equiv


0$). Поэтому из теоремы 1 следует также сходимость метода подобластей для


интегральных, дифференциальных и интегродифференциальных уравнений с


периодическими коэффициентами.


\end{note}








В частном случае теорема 1 несколько упрощается и усиливается, что показывает








\begin{thm}


Пусть $y(s) \in L_2$, а оператор $B$ таков, что вполне


непрерывен оператор $T:W_2^{m-1} \to C_{2 \pi}$, где $m \geq 1$ и


$$


Tx \equiv \int_{0}^{s} (Bx)( \sigma) d \sigma -


\frac{s}{2 \pi} \int_{0}^{2 \pi}(Bx)


( \sigma) d \sigma, \quad x \in W_2^{m-1}.


$$


Если задача $(1)$--$(2)$ при $y(s)=0$ имеет в $L_2$ лишь тривиальное решение,


то метод подобластей $(1)$--$(6)$ сходится в пространстве $ W_2^{m-1} $ со скоростью


\begin{equation}


\| x^* - x_n^* \|_{W_2^{m-1}}=O \lbrace E_n^T(x^{*(m-1)})_{C} \rbrace.


\end{equation}


\end{thm}








\begin{pf}


Следуя [9], [10], можно показать, что


задача (1)--(2) и СЛАУ (6) эквивалентны уравнениям соответственно


\begin{gather}


Kx \equiv x^{(m-1)}(s) - x^{(m-1)}(0) + \int_{0}^{s} (Bx)( \sigma) d \sigma -


\frac{s}{2 \pi} \int_{0}^{2 \pi} (Bx)( \sigma) d \sigma=\ov{y}(s), \\


%


K_n x_n \equiv x_n^{(m-1)}(s) - x_n^{(m-1)}(0) + \rL_n \bigg\lbrace \int_{0}^{s}


(Bx_{n})( \sigma) d \sigma -


\frac{s}{2 \pi} \int_{0}^{2 \pi} (Bx_n)( \sigma) d \sigma \bigg\rbrace=


\ov y_n(s),


\end{gather}


где


$$


\ov y(s)=\int_{0}^{s} y( \sigma) d \sigma - \frac{s}{2 \pi} \int_{0}^


{2 \pi} y( \sigma) d \sigma, \quad \ov y_n(s)=\rL_n \ov y(s),


$$


а $\rL_n$ --- определенный выше оператор Лагранжа по узлам (5).





Уравнение (27) будем рассматривать как операторное в пространстве


$$


\ov X=\lbrace x(s) \in C_{2 \pi}^{m-1} : x^{(k)}(0)=x^{(k)}


(2 \pi), \ \; k=\overline{0,m-1} \rbrace


$$


с нормой


$$


{\| x \|}_{\ov X}=\sum_{k=0}^{m-1} {\| x^{(k)}(s) \|}_{L_{2}}


=\| x \|_{W_2^{m-1}}, \quad x \in \ov X.


$$


Ясно, что в условиях теоремы $ K $ есть аддитивный и однородный оператор из


$\ov X $ в $ \ov Y=\lbrace y(s) \in C_{2 \pi} \rbrace $ с нормой $ L_2 $.


Пусть $ \ov X_n \subset \ov X $ и $ \ov Y_n \subset \ov Y $ ---


подпространства всех тригонометрических полиномов порядка не выше $ n $ с


нормами соответственно $ \ov X $ и $ \ov Y $. Тогда (28) можно рассматривать


как операторное уравнение в пространстве $ \ov X_n $; ясно, что $ K_n $


есть линейный оператор из $ \ov X_n $ в $ \ov Y_n $ при каждом


фиксированном натуральном $ n $.





Далее важно подчеркнуть, что (28) есть уравнение метода коллокации


для уравнения (27). С учетом этого факта в дальнейшем будем пользоваться


способом исследования метода коллокации, предложенным в [11].





В силу (27) и (28) для любого $x_n \in \ov X_n$, $x_n \neq 0$, с учетом


аппроксимативных свойств оператора Лагранжа $\rL_n $ в $ L_2 $


(напр., [4], гл.\,3) находим


\begin{gather*}


{\| Kx_n - K_nx_n \|}_2={\| Tx_n - \rL_n Tx_n \|}


_2=\| x_n \|\, {\| Tz_n - \rL_n Tz_n \|}_2 \leq


2 \| x_n \|_{\ov X_n} E_n^T(Tz_n)_{C},\\


z_n=x_n / \| x_n \|, \ \ C=C_{2 \pi}.


\end{gather*}


Отсюда и из теоремы Джексона [5]--[7] в $ C_{2 \pi} $ следует


\begin{gather*}


{\| Kx_n - K_nx_n \|}_{\ov Y} \leq 6 {\| x_n \|}_{\ov X}


w(Tz_n; \tfrac{1}{n}) \leq 6 {\| x_n \|}_{\ov X} \beta_n, \quad x_n


\in \ov X_n, \\


%


\beta_n=\sup \lbrace w(z; \tfrac{1}{n}) : z \in TS(0,1) \subset \ov Y \rbrace,


\end{gather*}


где $w( \varphi; \delta)$ --- модуль непрерывности функции


$\varphi \in C_{2 \pi}$ с шагом $\delta \in (0, \pi)$,


а $ S(0,1) \subset \ov X $.





В силу условия теоремы множество $TS(0,1) $ является компактным в $C_{2 \pi} $,


а тогда из теоремы 3.1 книги [7] следует $\beta_n \to 0 $ при $ n \to


\infty$. Поэтому


\begin{equation}


\varepsilon_n \equiv \| K - K_{n} \|_{\ov X_n \to \ov Y}=O( \beta


_n) \to 0, \quad n \to \infty.


\end{equation}





Теперь, применяя теорему 7 ([4], гл.\,1) к уравнениям (27) и (28), в силу (29)


находим, что операторы $K_n : \ov X_n \to \ov Y_n$, $n \gea n_0$,


линейно обратимы (следовательно, система (6) однозначно разрешима), и обратные


операторы $K_n^{-1} $ ограничены по норме в совокупности:


\begin{equation}


\| K_n^{-1} \|=O(1).


\end{equation}





Для получения оценки погрешности воспользуемся теоремой 6 ([4], гл.\,1)


применительно к уравнениям (27) и (28):


\begin{multline}


\| x^* - x_n^* \|_{\ov X}=\| K^{-1} \ov y - K_n^


{-1} \ov y_n \|_{\ov X} \leq \| ( E - K_n^{-1} \rL_n K)


(x^* - \ov x_n) \| \leq \\


%


\leq \| x_n^* - \ov x_n \|_{\ov X} + \| K_n^{-1} \| \,


 \| \rL_n K(x^* - \ov x_n) \|_{\ov Y},


\end{multline}


где $\ov x_n$ --- произвольный элемент из $\ov X_n$. Так как


$\|\rL_n \|_{C \to L_2}=1 $ (напр., [4], гл.\,3), то в силу (30),


(31) и (27) последовательно находим


\begin{multline}


\| x^* - x_n^* \|_{W_2^{m-1}} \leq \| x^* - \ov x_n \|_{W_


2^{m-1}} + d_1 \| K(x^* - \ov x_n) \|_{C} \leq \\


%


\leq \| x^* - \ov x_n \|_{W_2^{m-1}} + d_1 \bigg\lbrace \| [x^{*(m-1)}


(s) - \ov x_n^{(m-1)}(s)] - \\


%


- [x^{*(m-1)}(0) - \ov x_n^{(m-1)}(0)] \|_{C} + \bigg\| \int_{0}^{s}


B(x^* - \ov x_n) d \sigma - \frac{s}{2 \pi} \int_{0}^{2 \pi} B(x^* - \ov


x_n) d \sigma \bigg\|_{C} \bigg\rbrace \leq \\


%


\leq \| x^* - \ov x_n \|_{W_2^{m-1}} + d_1 \lbrace 2 \| x^{*(m-1)}


(s) - \ov x_n^{(m-1)}(s) \|_{C} + \| T(x^* - \ov x_n) \|_{C}


\rbrace \leq \\


%


\leq (1 + d_1d_2) \| x^* - \ov x_n \|_{W_2^{m-1}} +2d_1 \| x^{*


(m-1)} - \ov x_n^{(m-1)} \|_{C}.


\end{multline}





Поскольку


$$


\| x^* - \ov x_n \|_{W_2^{m-1}} \leq d_3 \| x^{*(m-1)} -


\ov x_n^{(m-1)} \|_{C},


$$


то в силу (32)


\begin{equation}


\| x^* - x_n^* \|_{W_2^{m-1}} \leq d_4 \| x^{*(m-1)} - \ov x_n^{


(m-1)} \|_{C}.


\end{equation}


Выбирая полином $ \ov x_n(s) $ так, чтобы $ \ov x_n^{(m-1)}(s) $ был


полиномом наилучшего равномерного приближения для $ x^{*(m-1)}(s) \in C $,


из (33) находим (26).


\end{pf}





Из теоремы 2 выводятся следующие утверждения.





\begin{cor}


В условиях теоремы 2 метод (1)--(6) сходится в пространствах


$C^k$, $k=\overline{0,m-1}$, со скоростями соответственно


\begin{align}


\| x^* - x_n^* \|_{C^k} &=O \lbrace E_n^T


(x^{*(m-1)})_{C} \rbrace, \quad k=\overline{0,m-2}; \\


%


\| x^* - x_n^* \|_{C^{m-1}} &=O \lbrace E_n^T


(x^{*(m-1)})_{C} \| \rL_n \|_{C} \rbrace=


O \lbrace E_n^T(x^{*(m-1)})_{C} \ln n \rbrace.


\end{align}


Пусть, кроме того, $y(s) \in C $, а оператор $ B:Z \to C$


ограничен, где $Z$ --- одно из пространств $W_2^{m-1}$,


$C^{m-1}$; тогда в пространстве $ C^m $ справедливы оценки


\begin{equation}


\| x^* - x_n^* \|_{C^m}=O \lbrace \|


P_n \|_{C} (E_n^T(x^{*(m)})_{C} + \|


x^* - x_n^* \|_Z) \rbrace=


O \lbrace E_n^T (x^{*(m)})_{C} \ln n \rbrace.


\end{equation}


\end{cor}








\begin{cor}


Пусть оператор $B$ и функция $y(s) $ таковы, что


$x^{*(m)}(s) \in H_{p}^{r + \alpha} $, где $r \geq 0$ целое, $0 < \alpha


\leq 1$, $1 \leq p \leq \infty$. Тогда в условиях теоремы 2 метод (1)--(6)


сходится в пространствах $W_2^m$, $C^k$ ($k=\overline{0,m}$) со скоростями


соответственно


\begin{alignat*}{2}


\| x^* - x_n^* \|_{W_2^{m-1}} &=O \bigg\lbrace \frac{1}{n^{r +


\alpha + 1/q}} \bigg\rbrace, &\quad r + \alpha + 1/q &> 0; \\


%


\| x^* - x_n^* \|_{C^k} &=O \bigg\lbrace \frac{1}{n^{r + \alpha


 + 1/q}} \bigg\rbrace, &\quad r + \alpha + 1/q &> 0 ,\ \ k=\overline{0,m-2}; \\


%


\| x^* - x_n^* \|_{C^{m-1}} &=O \bigg\lbrace \frac{


\ln n}{n^{r + \alpha + 1/q}} \bigg\rbrace, &\quad r + \alpha + 1/q &> 0; \\


%


\| x^* - x_n^* \|_{C^m} &=O \bigg\lbrace \frac{\ln


 n}{n^{r + \alpha - 1/p}} \bigg\rbrace, &\quad r + \alpha - 1/p &> 0,


\end{alignat*}


где $q$ --- сопряженный с $p$ показатель: $ 1/p + 1/q=1$,


$1 \leq p \leq \infty $.


\end{cor}








{\bf Доказательство следствия 3.}


Оценка (34) есть следствие неравенств (26) и (15). Докажем (35) при


условии $x^{*(m-1)}(0)=x_n^{*(m-1)}(0) $. Положим $ \rho_n(s)=x^{*(m-1)}(s) -


x_n^{*(m-1)}(s) $. Тогда из (27) при $ x=x^*(s) $ и из (28) при $ x_n=x_n^*(s)$


в узлах (5) находим


$$


\rho_n(s_j)=- \int_0^{s_j} B(x^* - x_n^*, \sigma) d \sigma + \frac{s_j}{2 \pi}


\int_0^{2 \pi} B(x^* - x_n^*, \sigma) d \sigma.


$$


Поэтому


$$


\max_{0 \leq j \leq 2n} | \rho_n(s_j) |=\max_{0 \leq j \leq 2n}


| T(x^* - x_n^*;s_j) | \leq \| T(x^* - x_n^*) \|_{C}


\leq d_2 \| x^* - x_n^* \|_{W_2^{m-1}}.


$$


Тогда


\begin{multline}


\| \rho_n(s) \|_{C}=\| x^{*(m-1)}(s) - x_n^{*(m-1)}(s) \|_{C}


\leq \| x^{*(m-1)}(s) - \rL_n x^{*(m-1)}(s) \|_{C} + \\


%


+ \| \rL_n \rho_n(s) \|_{C} \leq 2 \| \rL_n \|_{C} 


E_n^T(x^{*(m-1)})_{C} + \| \rL_n \|_{C} \max_{0 \leq j \leq 2n}


| \rho_n(s_j) |=\\


%


=O \lbrace \| \rL_n \|_{C} (E_n^T(x^{*(m-1)})_{C}


+ \| x^* - x_n^* \|_{W_2^{m-1}}) \rbrace.


\end{multline}





Из (37) и (26) с учетом известной оценки $ \| \rL_n \|_C=O( \ln n) $ находим


\begin{equation}


\| \rho_n(s) \|_{C}=\| x^{*(m-1)} - x_n^{*(m-1)} \|_{C}=


O \lbrace E_n^T(x^{*(m-1)})_{C} \ln n \rbrace.


\end{equation}


Теперь из (38) и (34) получаем оценку (35)


$$


\| x^* - x_n^* \|_{C^{m-1}}=\| x^* - x_n^* \|_{C^{m-2}} +


\| x^{*(m-1)} - x_n^{*(m-1)} \|_{C}=O \lbrace E_n^T(x^{*(m-1)})_


{C} \ln n \rbrace.


$$





Для доказательства (36) воспользуемся тождествами


$$


x^{*(m)} \equiv y - Bx^*, \quad x_n^{*(m)} \equiv P_ny - P_nBx_n^*.


$$


Отсюда последовательно находим


\begin{multline}


\| x^{*(m)} - x_n^{*(m)} \|_{C} \leq \| x^{*(m)} - P_nx^{*(m)}


\|_{C} + \| P_nB(x^* - x_n^*) \|_{C} \leq \\


%


\leq 2 \| P_n \|_{C} E_n^T(x^{*(m)})_{C} + \| P_n \|_{C}


\| B(x^* - x_n^*) \|_{C} \leq \\


%


\leq \| P_n \|_{C} \lbrace 2E_n^T(x^{*(m)})_{C} + \| B \|_


{Z \to C} \| x^* - x_n^* \|_{Z} \rbrace=\\


%


=O \lbrace \| P_n \|_{C \to C} ( E_n^T(x^{*(m)})_{C} +


\| x^* - x_n^* \|_{Z} ) \rbrace.


\end{multline}





В силу (26) и (35)


\begin{align}


\| x^* - x_n^* \|_{W_2^{m-1}} &=O \lbrace E_n^T(x^{*(m-1)})_{C}


\rbrace=O \bigg\lbrace \frac{E_n^T(x^{*(m)})_{C}}{n} \bigg\rbrace, \\


%


\| x^* - x_n^* \|_{C^{m-1}} &=O \bigg\lbrace E_n^T(x^{*(m)})_{C}


\frac{\ln n}{n} \bigg\rbrace.


\end{align}


Поэтому $\| x^* - x_n^* \|_{Z}=o \lbrace E_n^T(x^{*(m)})_{C}


\rbrace $. Тогда из (39) с учетом (24), (25) получаем оценку


$$


\| x^{*(m)} - x_n^{*(m)} \|_{C}=O \lbrace \| P_n \|_{C}


\ E_n^T(x^{*(m)})_{C} \rbrace=O \lbrace E_n^T(x^{*(m)})_{C} \ln n


\rbrace.


$$


Отсюда и из (39)--(41) следует оценка (36).$\quad\square$


\medskip





Следствие 4 доказывается


аналогично следствию 2 теоремы 1 с применением теорем типа Джексона в


$C_{2 \pi} $ (напр., [5]--[8]).





\begin{note}


В теореме 2 в качестве оператора $B$ можно взять


интегродифференциальный оператор $(1')$ лишь при $h_m(s, \sigma) \equiv 0$;


в этом и состоит недостаток теоремы 2 по сравнению с теоремой 1. В то же время


при выполнении этого дополнительного условия теорема 2 имеет и ряд преимуществ


перед теоремой 1.


\end{note}





Следует также отметить, что при выборе в качестве $B$ дифференциального


оператора


$$


Bx \equiv \sum_{k=0}^{m-1} a_k(s) x^{(k)}(s),


\quad m \geq 1,\ \ a_k(s) \in L_p,\ \ p >1,


$$


где правая часть $(1)$ --- функция $ y(s) \in L_p $ при $ p>1 $, менее общий


результат, чем теорема 2, но другим (на наш взгляд, более сложным) способом


получен в [9], [10].





\begin{thebibliography}{99}





\bibitem{1}


Габдулхаев Б.Г., Ермолаева Л.Б. {\em Один новый полиномиальный оператор и его


приложения} // Теория прибл. функций. -- М.: Наука, 1987. -- С.\,98--100.


\bibitem{2}


Ермолаева Л.Б. {\em Аппроксимативные свойства полиномиальных операторов и решение


интегральных и интегродифференциальных уравнений методом подобластей}.


-- Дис.~$\dotsc$ канд. физ.-матем. наук. -- Казань, 1987. -- 154~с.


\bibitem{3}


Канторович Л.В., Акилов Т.П. {\em Функциональный анализ в нормированных


пространствах}. -- М.: Наука, 1977. -- 684~с.


\bibitem{4}


Габдулхаев Б.Г. {\em Оптимальные аппроксимации решений линейных задач}. --


Казань: Изд-во Казанск. ун-та, 1980. -- 232~с.


\bibitem{5}


Ахиезер Н.И. {\em Лекции по теории аппроксимации}. -- М.: Наука, 1965. -- 407~с.


\bibitem{6}


Даугавет И.К. {\em Введение в теорию приближения функций}. -- Л.: Изд-во


ЛГУ, 1977. -- 184~с.


\bibitem{7}


Тиман А.Ф. {\em Теория приближения функций действительного переменного}. --


М.: Физматгиз, 1960. -- 624~с.


\bibitem{8}


Андриенко В.А. {\em Теоремы вложения для функций одного переменного} // Итоги науки


и техники. Матем. анализ. -- М.: ВИНИТИ, 1971. -- С.\,203--262.


\bibitem{9}


Керге Р.М. {\em Метод подобластей для периодической задачи} // Учен. зап. Тартуск.


ун-та. -- 1979. -- \No\,500. -- С.\,53--68.


\bibitem{10}


Керге Р.М. {\em О сходимости и устойчивости метода подобластей}: Автореф.


дис.~$\dotsc$ канд. физ.-матем. наук. -- Тарту, 1979. -- 10~с.


\bibitem{11}


Лучка А.Ю., Каспшицкая М.Ф. {\em О методе коллокации} // Журн. вычисл. матем. и


матем. физ. -- 1968. -- Т.\,8. -- \No\,5. -- С.\,950--964.





\end{thebibliography}





\vskip 2cm





\post{\it Казанская государственная}{\it Поступила}


\post{\it архитектурно-строительная академия}{11.03.2003}





\label{end}


\end{document}











%\underline{Лемма 1}. Операторы $ P_n=P_n^2 : L_2 \to L_2 $ сильно сходятся


%к единичному оператору $ E $ пространства $ L_2 $, причем для любых $ n \in N $


%справедливы неравенства








%\underline{Следствие 1}. Справедливы следующие порядковые оценки





%\underline{Следствие 2}. Последовательность операторов $ \lbrace P_n \rbrace $


%не может оказаться сходящейся на всем пространстве $ С $. В то же время для


%любой функции $ \varphi(s) \in C $, удовлетворяющей условию Дини--Липшица,


%полиномы $ P_n( \varphi;s) $ равномерно сходятся со скоростью


%


%$$


%\| \varphi - P_n \varphi \|_{C} \leq 2 \| P_n \|_{C} E_n^T


%( \varphi)_{C}=O \lbrace E_n^T( \varphi)_{C} \ln n \rbrace.


%$$


%


%\hskip 0.5mm





Для завершения доказательства остается показать справедливость неравенств


(15). Очевидно, что в доказательстве нуждаются лишь вторые части неравенств.


Пусть сначала $k=\overline{1,m-1} $. Рассмотрим краевую задачу


$$


x''(s)=x''(s), \quad x(0)=x(2 \pi), \quad x'(0)=x'(2 \pi).


$$





По теореме Ролля существует точка $b_1 \in (0,2 \pi)$ такая, что


$x'(b_{1})=0$. Поэтому имеем тождество


\begin{equation}


x'(s)=\int _{b_1}^{s} x''( \sigma) d \sigma,


\tag{$\alpha$}


\end{equation}


откуда и следует (15) при $k=1$, $m=2$, $d_0=2 \pi$.





По аналогии с ($\alpha$) получаем тождество


\begin{equation}


x''(s)=\int_{b_2}^{s} x'''( \sigma) d \sigma,\quad x \in W_2^3,


\tag{$\beta$}


\end{equation}


где $b_2$ --- такая постоянная, что $x''(b_2)=0$; эта постоянная


существует в силу теоремы Ролля и условия $x'(0)=x'(2 \pi)$.


Из $(\beta)$ следует оценка (15) при $k=2$, $m=3$, $d_0=2 \pi$.





По аналогии с $(\alpha)$, $(\beta)$ находим тождество


\begin{equation}


x^{(i)}(s)=\int_{b_i}^s x^{(i+1)}( \sigma) d \sigma, \quad x \in W_2^{i+1},


\quad i={\overline{0,m-1}},


\tag{$\gamma$}


\end{equation}


где $b_i \in (0, 2 \pi)$ --- такие постоянные, что $x^{(i)}(b_i)=0$.


Теперь из $(\gamma)$ следует (15) при всех $k=i$, $m=i+1$, $d_0=2 \pi$:


\begin{equation}


\| x^{(i)} \|_{L_2} \leq \| x^{(i)} \|_{C} \leq 2 \pi \|


x^{(i+1)} \|_{L_2},\ \ i={\overline{1,m-1}}.


\tag{$\delta$}


\end{equation}





Повторное применение неравенств $(\delta)$ приводит к (15) при 


$d_0=(2 \pi)^{m-1}$, $k={\overline{1,m-1}}$.





Пусть теперь $k=0$. Рассмотрим краевую задачу


\begin{equation}


x'(s)=x'(s),\ \ x(0)=x(2 \pi).


\tag{$\ve$}


\end{equation}


Тогда для любой $s \in [0,2 \pi]$ имеем


\begin{equation}


x(s)=\int_{a}^{s} x'( \sigma) d \sigma + x(a).


\tag{$\mu$}


\end{equation}


Постоянную $a$ выберем так, чтобы


$$


x(a)=\frac{1}{2 \pi} \int_{0}^{2 \pi} x( \sigma) d \sigma;


$$


поскольку $x(s) \in C_{2 \pi}=C$, то по теореме о среднем интегрального


исчисления такая постоянная всегда существует. Поэтому справедливо тождество


\begin{equation}


x(s)=\int_{a}^{s} x'( \sigma) d \sigma + \frac{1}{2 \pi} \int_{0}^{2 \pi}


x( \sigma) d \sigma.


\tag{$\mu'$}


\end{equation}


Отсюда следуют неравенства


$$


\| x \|_{C} \leq 2 \pi \| x' \|_{L_2}


+ \| x \|_{L_2} \leq 2 \pi \| x \|_{W_2^1}.


$$





Если существует хотя бы одна точка $a \in [0, 2 \pi]$ такая, что $x(a)=0$,


то из соотношений $(\mu)$, $(\mu')$ следует неравенство (15) при $m=1$,


$k=0$. Если же такой точки нет, то вместо $(\ve)$ рассматривается краевая


задача


$$


x'(s) + px(s)=x'(s) + px(s), \quad x(0)=x(2 \pi),


$$


где $p$ --- такая постоянная, что задача $x'+px=0$, $x(0)=x(2 \pi)$ имеет


лишь нулевое решение.


\end{pf}


